

% \label{sec:rq2}

% % Having identified Gemini-3-Pro as the strongest verifier in RQ1, we employ it to evaluate specifications generated by all models on \textsc{Coins}, measuring their \textsc{Reject}$_{\text{all}}$ rates.

% Having identified Gemini 3 Pro Preview
% as the strongest verifier in RQ1, we employ it as the primary verifier to evaluate specifications generated by all six models on \textsc{Coins}. Table~\ref{tab:generation} presents the results, where \textsc{Reject}$_{\text{all}}$ serves as our core metric for identifying candidate specifications.


% \begin{table*}[htbp]
% \centering
% \setlength{\aboverulesep}{0pt}
% \setlength{\belowrulesep}{0pt}
% \caption{Specification Generation Results (n=164). {\footnotesize\textit{Syntax denotes syntactically valid specifications. Higher values indicate better performance.}}}
% \label{tab:generation}
% \small
% \newcolumntype{Y}{>{\columncolor{yellow!15}}c}
% \begin{tabular}{@{} cc ccc Y @{}}
% \toprule
% \textbf{Generate Model} & \textbf{Verify Model} & \textbf{Syntax} & \textbf{\textsc{Pass}$_{\text{first}}$} & \textbf{\textsc{Pass}$_{\text{all}}$} & \cellcolor{yellow!15}\textbf{\textsc{Reject}$_{\text{all}}$} \\
% \midrule
% \multirow{2}{*}{GPT-4o} 
%   & Gemini 3 Pro Preview & \multirow{2}{*}{14.63\% \textcolor{gray}{(24)}} & 12.20\% \textcolor{gray}{(20)} & 4.27\% \textcolor{gray}{(7)} & 4.27\% \textcolor{gray}{(7)} \\
%   & GPT-4o &  & 0.61\% \textcolor{gray}{(1)} & 0.61\% \textcolor{gray}{(1)} & -- \\
% \midrule
% \multirow{2}{*}{DeepSeek-V3.1} 
%   & Gemini 3 Pro Preview & \multirow{2}{*}{7.93\% \textcolor{gray}{(13)}} & 4.27\% \textcolor{gray}{(7)} & 1.22\% \textcolor{gray}{(2)} & 1.22\% \textcolor{gray}{(2)} \\
%   & DeepSeek-V3.1 &  & 1.22\% \textcolor{gray}{(2)} & 0.00\% \textcolor{gray}{(0)} & -- \\
% \midrule
% \multirow{2}{*}{Claude 3.7 Sonnet} 
%   & Gemini 3 Pro Preview & \multirow{2}{*}{19.51\% \textcolor{gray}{(32)}} & 15.24\% \textcolor{gray}{(25)} & 1.83\% \textcolor{gray}{(3)} & 1.83\% \textcolor{gray}{(3)} \\
%   & Claude 3.7 Sonnet &  & 4.27\% \textcolor{gray}{(7)}  & 0.00\% \textcolor{gray}{(0)} & -- \\
% \midrule
% \multirow{2}{*}{Claude 4.5 Opus} 
%   & Gemini 3 Pro Preview & \multirow{2}{*}{59.76\% \textcolor{gray}{(98)}} & 45.12\% \textcolor{gray}{(74)} & 14.63\% \textcolor{gray}{(24)} & 14.63\% \textcolor{gray}{(24)} \\
%   & Claude 4.5 Opus &  & 50.61\% \textcolor{gray}{(83)} & 10.98\% \textcolor{gray}{(18)} & -- \\
% \midrule
% \multirow{2}{*}{GPT-5} 
%   & Gemini 3 Pro Preview & \multirow{2}{*}{67.07\% \textcolor{gray}{(110)}} & 60.98\% \textcolor{gray}{(100)} & 15.24\% \textcolor{gray}{(25)} & 14.63\% \textcolor{gray}{(24)} \\
%   & GPT-5 &  & 28.66\% \textcolor{gray}{(47)} & 7.32\% \textcolor{gray}{(12)} & -- \\
% \midrule
% Gemini 3 Pro Preview & Gemini 3 Pro Preview & 78.05\% \textcolor{gray}{(128)} & 60.98\% \textcolor{gray}{(100)} & 28.05\% \textcolor{gray}{(46)} & 26.22\% \textcolor{gray}{(43)} \\
% \bottomrule
% \end{tabular}
% \end{table*}

% % \textbf{Generation capability mirrors verification capability.} The results reveal a notable consistency: models that excel at verification also tend to generate higher-quality specifications. Gemini-3-Pro again achieves the best performance, suggesting that specification generation and proof synthesis may share underlying reasoning capabilities.


% \paragraph{Generation Capability Mirrors Verification Capability.}
% The results reveal a striking correlation: models that excel at verification also produce higher-quality specifications. Gemini-3-Pro achieves the best overall performance with 78.05\% syntax correctness and 26.22\% \textsc{Reject}$_{\text{all}}$, followed by GPT-5 (67.07\%, 15.24\%) and Claude-4.5-Opus (59.76\%, 14.63\%). This consistency is unsurprising, as both tasks fundamentally require generating valid Rocq code—specifications are Rocq definitions, and proofs are Rocq tactics. Models with stronger Rocq fluency naturally excel at both.


% % \textbf{Substantial performance gap.} We observe significant disparities across models. In terms of syntactic correctness, GPT-5, Gemini-3-Pro, and Claude-Opus-4.5 significantly outperform GPT-4o, DeepSeek-V3.1, and Claude-Sonnet-3.7. This gap persists in the \textsc{Pass}$_1$ and \textsc{Pass}$_{\text{all}}$ stages, confirming that generating verifiable formal specifications remains a challenging task.

% \paragraph{Substantial Performance Gap.}
% We observe significant disparities between frontier and earlier-generation models. The syntax correctness gap is pronounced: Gemini-3-Pro (78.05\%), GPT-5 (67.07\%), and Claude-4.5-Opus (59.76\%) substantially outperform Claude-3.7-Sonnet (19.51\%), GPT-4o (14.63\%), and DeepSeek-V3.1 (7.93\%). This disparity amplifies at subsequent stages—the best earlier-generation model (GPT-4o) achieves only 4.27\% \textsc{Reject}$_{\text{all}}$, compared to 26.22\% for Gemini 3 Pro Preview. These results confirm that generating formally verifiable specifications remains challenging and benefits significantly from recent advances in reasoning capabilities.


% % \textbf{Limited self-verification bias.} To rule out potential bias from a single external verifier, we additionally evaluate each model's specifications under self-verification. The differences are marginal for strong models, indicating that the choice of verifier minimally affects results when the verifier is sufficiently capable. However, weaker verifiers consistently underestimate specification quality—they fail to prove correct specifications, resulting in deflated \textsc{Pass}$_{\text{all}}$ scores.


% % \textbf{Limited self-verification bias.}
% % To rule out potential bias from a single verifier, we additionally evaluate each model's specifications under self-verification. For strong models, differences are marginal: Claude-4.5-Opus achieves 10.98\% \textsc{Pass}$_{\text{all}}$ under self-verification versus 14.63\% under Gemini-3-Pro. However, weaker verifiers systematically underestimate specification quality by failing to prove correct specifications. Notably, Claude-3.7-Sonnet produces no valid proofs under self-verification.


% \paragraph{Limited Self-Verification Bias.}
% A natural concern is whether Gemini-3-Pro's leading performance benefits from evaluating its own generated specifications. To investigate this potential bias, we conduct self-verification experiments where each model verifies its own outputs. The results indicate that self-verification does not inflate scores: strong models achieve comparable or lower \textsc{Pass}$_{\text{all}}$ under self-verification than under Gemini-3-Pro verification (e.g., Claude-4.5-Opus: 10.98\% vs.\ 14.63\%; GPT-5: 7.32\% vs.\ 15.85\%). Conversely, weaker verifiers systematically underestimate specification quality due to their limited ability to construct valid proofs.


% % \textbf{Verification remains challenging even for test acceptance.} Even proving that a specification accepts the provided test cases—a relatively basic verification objective—proves difficult for many models. This underscores a limitation of benchmarks with complex verification tasks:


% \paragraph{Verification Remains Challenging even for Test Acceptance.}
% Even proving that a specification accepts the provided test cases—a relatively basic verification objective—proves difficult for many models. The drop from \textsc{Pass}$_1$ to \textsc{Pass}$_{\text{all}}$ is substantial across all models: Gemini-3-Pro drops from 60.98\% to 28.05\%, GPT-5 from 60.98\% to 15.24\%, and Claude-4.5-Opus from 45.12\% to 14.63\%. This reveals that while models can often prove acceptance for individual test cases, achieving consistent verification across all cases remains elusive. The gap underscores a fundamental limitation of benchmarks with complex verification tasks:



% \vspace{0.5em}
% \label{trade-off}
% \begin{tcolorbox}[
%   colback=gray!10,
%   colframe=gray!60,
%   boxrule=0.6pt,
%   arc=2pt,
%   left=6pt,
%   right=6pt,
%   top=6pt,
%   bottom=6pt
% ]
% \small
% \textbf{Verification Complexity Trade-off.} The harder the verification task, the stronger the prover required, and the more difficult it becomes to accurately assess specification quality. Even when the verification objective perfectly captures specification correctness, overly complex proof obligations risk systematic underestimation, rendering such benchmarks unreliable metrics for specification generation.
% \end{tcolorbox}
% \vspace{0.5em}




\label{sec:rq2}

Having identified Gemini 3 Pro Preview
as the strongest verifier in RQ1 (Sec~\ref{sec:rq1}), we employ it as the primary verifier to evaluate specifications generated by all six models on \textsc{Coins}. Table~\ref{tab:generation} presents the results, where \textsc{Reject}$_{\text{all}}$ serves as our core metric for identifying candidate specifications.

\textbf{Syntax Correctness Is Still Challenging.}
Generating syntactically valid Rocq specifications remains a fundamental challenge for most models. While Gemini 3 Pro Preview achieves 78.05\% syntax correctness and GPT-5 reaches 67.07\%, earlier-generation models struggle significantly—DeepSeek-V3.1 produces only 7.93\% valid specifications, and GPT-4o achieves just 14.63\%. This syntactic barrier filters out the majority of attempts before semantic evaluation can even begin. We provide a detailed analysis of syntax error causes in Appendix~\ref{appendix:rq7}.

% \textbf{Generation Capability Mirrors Verification Capability.}
% Models excelling at verification also produce higher-quality specifications. Gemini-3-Pro achieves the best performance with 78.05\% syntax correctness and 26.22\% \textsc{Reject}$_{\text{all}}$, followed by GPT-5 (67.07\%, 14.63\%) and Claude-4.5-Opus (59.76\%, 14.63\%). This consistency is expected, as both tasks require generating valid Rocq code.

\textbf{Substantial Performance Gap.}
Significant disparities exist between frontier and earlier-generation models. Syntax correctness ranges from 78.05\% (Gemini 3 Pro Preview) to 7.93\% (DeepSeek-V3.1), with this gap amplifying at later stages—GPT-4o achieves only 4.27\% \textsc{Reject}$_{\text{all}}$ versus 28.05\% for Gemini 3 Pro Preview.

\textbf{Limited Self-Verification Bias.}
Self-verification experiments confirm Gemini 3 Pro Preview's lead is not due to evaluating its own outputs: strong models achieve comparable or lower \textsc{Pass}$_{\text{all}}$ under self-verification (e.g., Claude 4.5 Opus: 10.98\% vs.\ 14.63\%).

\textbf{Precision Emerges from Coverage.}
The close alignment between \textsc{Reject}$_{\text{all}}$ and \textsc{Pass}$_{\text{all}}$ indicates that specifications passing all test cases are typically precise enough to reject incorrect implementations. This can be attributed to two factors: our diverse test suite effectively constrains the semantic space, and for relatively simple programs, there is little room for over-approximation in specification semantics, which aligns with the findings in~\cite{le2025can}.

\textbf{Comprehensive Test Suites Are Essential.}
The sharp drop from \textsc{Pass}$_{\text{first}}$ to \textsc{Pass}$_{\text{all}}$ across all models (e.g., Gemini 3 Pro Preview: 60.98\%$\to$28.05\%) confirms that individual test cases are insufficient—multiple tests collectively constrain specification semantics and distinguish genuine correctness from accidental success.


\begin{tcolorbox}[
  colback=gray!10,
  colframe=gray!60,
  boxrule=0.6pt,
  arc=2pt,
  left=6pt,
  right=6pt,
  top=6pt,
  bottom=6pt
]
\small
\textbf{Syntax Is the First Bottleneck.} For relatively simple algorithmic tasks such as HumanEval, syntactically valid Rocq specifications are often semantically precise under our test-based evaluation. The dominant failure mode is not that models systematically choose the wrong behavioral intent, but that they hallucinate invalid Rocq syntax or ill-typed formal constructs while trying to express program behavior.
\end{tcolorbox}


\begin{table*}[htbp]
\centering
\setlength{\aboverulesep}{0pt}
\setlength{\belowrulesep}{0pt}
\caption{Specification Generation Results (n=164). {\footnotesize\textit{Syntax denotes syntactically valid specifications. Higher values indicate better performance.}}}
\label{tab:generation}
\small
\newcolumntype{Y}{>{\columncolor{yellow!15}}c}
\resizebox{0.85\linewidth}{!}{
\begin{tabular}{@{} cc ccc Y @{}}
\toprule
\textbf{Generate Model} & \textbf{Verify Model} & \textbf{Syntax} & \textbf{\textsc{Pass}$_{\text{first}}$} & \textbf{\textsc{Pass}$_{\text{all}}$} & \cellcolor{yellow!15}\textbf{\textsc{Reject}$_{\text{all}}$} \\
\midrule
\multirow{2}{*}{GPT-4o} 
  & Gemini 3 Pro Preview & \multirow{2}{*}{14.63\% \textcolor{gray}{(24)}} & 12.20\% \textcolor{gray}{(20)} & 4.27\% \textcolor{gray}{(7)} & 4.27\% \textcolor{gray}{(7)} \\
  & GPT-4o &  & 0.61\% \textcolor{gray}{(1)} & 0.61\% \textcolor{gray}{(1)} & -- \\
\midrule
\multirow{2}{*}{DeepSeek-V3.1} 
  & Gemini 3 Pro Preview & \multirow{2}{*}{7.93\% \textcolor{gray}{(13)}} & 4.27\% \textcolor{gray}{(7)} & 1.22\% \textcolor{gray}{(2)} & 1.22\% \textcolor{gray}{(2)} \\
  & DeepSeek-V3.1 &  & 1.22\% \textcolor{gray}{(2)} & 0.00\% \textcolor{gray}{(0)} & -- \\
\midrule
\multirow{2}{*}{Claude 3.7 Sonnet} 
  & Gemini 3 Pro Preview & \multirow{2}{*}{19.51\% \textcolor{gray}{(32)}} & 15.24\% \textcolor{gray}{(25)} & 1.83\% \textcolor{gray}{(3)} & 1.83\% \textcolor{gray}{(3)} \\
  & Claude 3.7 Sonnet &  & 4.27\% \textcolor{gray}{(7)}  & 0.00\% \textcolor{gray}{(0)} & -- \\
\midrule
\multirow{2}{*}{Claude 4.5 Opus} 
  & Gemini 3 Pro Preview & \multirow{2}{*}{59.76\% \textcolor{gray}{(98)}} & 45.12\% \textcolor{gray}{(74)} & 14.63\% \textcolor{gray}{(24)} & 14.63\% \textcolor{gray}{(24)} \\
  & Claude 4.5 Opus &  & 50.61\% \textcolor{gray}{(83)} & 10.98\% \textcolor{gray}{(18)} & -- \\
\midrule
\multirow{2}{*}{GPT-5} 
  & Gemini 3 Pro Preview & \multirow{2}{*}{67.07\% \textcolor{gray}{(110)}} & 46.34\% \textcolor{gray}{(76)} & 15.24\% \textcolor{gray}{(25)} & 15.24\% \textcolor{gray}{(25)} \\
  & GPT-5 &  & 28.66\% \textcolor{gray}{(47)} & 7.32\% \textcolor{gray}{(12)} & -- \\
\midrule
Gemini 3 Pro Preview & Gemini 3 Pro Preview & 78.05\% \textcolor{gray}{(128)} & 60.98\% \textcolor{gray}{(100)} & 28.05\% \textcolor{gray}{(46)} & 28.05\% \textcolor{gray}{(46)} \\
\bottomrule
\end{tabular}
}
\end{table*}



% \textbf{Verification Remains Challenging even for Test Acceptance.}
% The drop from \textsc{Pass}$_\text{first}$ to \textsc{Pass}$_{\text{all}}$ is substantial: Gemini-3-Pro (60.98\%$\to$28.05\%), GPT-5 (60.98\%$\to$15.24\%), and Claude-4.5-Opus (45.12\%$\to$14.63\%). While models can often verify individual test cases, consistent verification across all cases remains elusive.The gap underscores a fundamental limitation of benchmarks with complex verification tasks:

% \begin{tcolorbox}[
%   colback=gray!10,
%   colframe=gray!60,
%   boxrule=0.6pt,
%   arc=2pt,
%   left=6pt,
%   right=6pt,
%   top=6pt,
%   bottom=6pt
% ]
% \small
% \label{trade-off}
% \textbf{Verification Complexity Trade-off.} The harder the verification task, the stronger the prover required, and the more difficult it becomes to accurately assess specification quality. Even when the verification objective perfectly captures specification correctness, overly complex proof obligations risk systematic underestimation, rendering such benchmarks unreliable metrics for specification generation.
% \end{tcolorbox}
